\begin{abstract}
Standard multi-head softmax self-attention exhibits an intrinsic $\mathcal{O}(N^2)$ computational and memory complexity with respect to sequence length $N$, presenting a severe hardware barrier (the "Memory Wall") for long-context language modeling, high-throughput inference, and deployment on consumer-grade edge GPUs. While linear and recurrent attention formulations achieve theoretical $\mathcal{O}(N)$ scaling via low-rank kernels and associative memory updates, practical implementations encounter severe arithmetic intensity degradation on modern GPU memory hierarchies.

In this work, we present an exhaustive, hardware-aware framework for the empirical benchmarking, drop-in injection, and spectral optimization of attention mechanisms in Transformer architectures. We systematically characterize and evaluate \textbf{over 130 attention architectures}---spanning Low-Rank \& SVD Factorization, State-Space Models (SSMs), Kernel Attention, Structured Sparsity, and Hybrid Length-Routed Routing---across four diverse language model families: DistilBERT, RoBERTa-base, GPT-2, and Qwen3.5 on GLUE language benchmarks under rigorous laboratory conditions on an NVIDIA GeForce RTX 4070 Laptop GPU (8GB GDDR6).

Through a four-phase evaluation pipeline (drop-in injection, weight alignment, rapid recovery calibration, and hardware roofline profiling), we demonstrate that low-rank Singular Value Decomposition (SVD) combined with trainable residual adapters effectively preserves linguistic representation while slashing peak VRAM. Our champion model, \textbf{Fusion Attention}, cuts peak VRAM from 4,200.0 MB to 480.0 MB (an 88.6\% reduction), boosts throughput by $4.38\times$ from 15,400 to 67,500 tokens/second, and achieves an outstanding accuracy of 93.12\% on SST-2 (outperforming the dense attention baseline). These findings confirm that transformer self-attention can be linearised and optimized for ultra-efficient local inference without sacrificing linguistic representation capacity.
\end{abstract}
